\chapter{特例坍缩——经典有限维凸优化}

\label{chap:08}

前七章一直在为无穷维情形搭建工作台：第 \ref{sec:3-2} 节把紧性从范数拓扑转移到弱拓扑与弱$^\ast$拓扑，第 \ref{chap:04} 章和第 \ref{chap:05} 章把分离、次微分与下半连续性提升到泛函层面，第 \ref{chap:06} 章与第 \ref{chap:07} 章则把对偶、锥约束和 KKT 系统写成统一的原--对偶框架。本章的任务不是另起炉灶，而是反过来说明：当空间退回到有限维时，这些抽象工具为何几乎全部自动简化，并最终坍缩成 Boyd 书中的经典有限维凸优化母语。

这种坍缩沿四条主线展开。第 \ref{sec:8-1} 节先解释有限维范数等价、弱拓扑与强拓扑一致、闭有界紧等事实，从而说明存在性与正则性的许多技术门槛为何会在 $\mathbb R^n$ 中自动消失。第 \ref{sec:8-2} 节在此基础上把抽象锥规划和 KKT 系统退回到 LP/QP 的矩阵公式，说明线性规划与二次规划并不是抽象理论之外的例外，而是第 \ref{chap:07} 章结果的 Euclidean 版。第 \ref{sec:8-3} 节把第 \ref{sec:7-1} 节的锥语言与第 \ref{sec:3-4} 节的谱论语言压缩到 SOCP 与 SDP，展示“迹内积”“半正定锥”“特征值分解”如何共同承担有限维对偶配对。最后第 \ref{sec:8-4} 节讨论几何规划，说明它虽然外观上带有乘法结构，但在对数坐标下仍然是第 \ref{sec:6-4} 节 Fenchel--Rockafellar 对偶的一个经典有限维坍缩。

因此，这一章最重要的阅读态度不是“终于回到熟悉公式”，而是“识别哪些熟悉公式其实依赖了维数有限”。Boyd 之所以能把大量论证写成矩阵代数和几何图形，是因为有限维里弱紧性、内点、伴随与矩阵转置、对偶球紧性等结构都被压缩到了一个共同背景中。本章要做的，正是把这种背景重新显式化，然后再把它交回给经典有限维凸优化。这样，第 \ref{chap:09} 章和第 \ref{chap:10} 章中的算法才会被理解为“对抽象原理的计算实现”，而不是只对有限维模型成立的技巧。

\section{欧氏空间的退化}

\label{sec:8-1}

这一节是全章的释压阀。前面几章里那些关于弱拓扑、弱紧性、相对内部、Slater 条件和 KKT 乘子的谨慎区分，在 $\mathbb R^n$ 中会出现大规模坍缩：不同范数彼此等价，弱拓扑与通常拓扑一致，闭有界集变成紧集，很多资格条件也随之变得更容易验证。有限维凸优化之所以显得“天然顺手”，并不是因为它遵循另一套理论，而是因为抽象理论中的许多门槛在这里自动消失了。

\subsection{有限维范数与拓扑的一致性}

\begin{proposition}[有限维中的范数等价]\label{prop:finite-dimensional-norm-equivalence}
设 $X$ 为有限维实向量空间，$\|\cdot\|_a$ 与 $\|\cdot\|_b$ 是 $X$ 上两个范数。则存在常数 $c,C>0$，使得对任意 $x\in X$，
\[
c\|x\|_a\le \|x\|_b\le C\|x\|_a.
\]
\end{proposition}

\begin{proof}
在单位球
\[
S_a:=\{x\in X:\|x\|_a=1\}
\]
上考虑连续函数 $x\mapsto \|x\|_b$。由于 $S_a$ 在有限维空间中是紧集，故该函数取到正的最小值 $m$ 和有限的最大值 $M$。于是对任意非零 $x$，令 $u=x/\|x\|_a$，便有
\[
m\le \|u\|_b\le M.
\]
两边乘以 $\|x\|_a$ 即得
\[
m\|x\|_a\le \|x\|_b\le M\|x\|_a.
\]
令 $c=m$、$C=M$ 即可。
\end{proof}

\begin{proposition}[有限维中弱拓扑与强拓扑一致]\label{prop:finite-dimensional-weak-strong-equivalence}
设 $X$ 为有限维赋范空间，则第 \ref{sec:3-2} 节定义~\ref{def:weak-topology} 所给出的弱拓扑与范数拓扑一致。
\end{proposition}

\begin{proof}
取 $X$ 的一组基 $\{e_1,\dots,e_n\}$，并记其对偶基为 $\{e_1^\ast,\dots,e_n^\ast\}$。若网 $x_\alpha$ 在弱拓扑下收敛到 $x$，则对每个 $i$，
\[
e_i^\ast(x_\alpha)\to e_i^\ast(x).
\]
因此坐标逐项收敛。有限维中任意范数都与 Euclidean 范数等价，由命题~\ref{prop:finite-dimensional-norm-equivalence} 可知坐标逐项收敛推出范数收敛，所以弱收敛蕴含强收敛。反向蕴含则对任意空间都成立，因为每个连续线性泛函在范数拓扑下本就连续。故两种拓扑一致。
\end{proof}

\begin{corollary}[有限维中的闭有界紧性]\label{cor:finite-dimensional-compactness}
设 $X$ 为有限维赋范空间，则子集 $C\subset X$ 在下列三条中彼此等价：
\begin{enumerate}
  \item $C$ 闭且有界；
  \item $C$ 在范数拓扑下紧；
  \item $C$ 在弱拓扑下紧。
\end{enumerate}
\end{corollary}

\begin{proof}
由 Heine--Borel 定理，闭且有界等价于范数紧。再由命题~\ref{prop:finite-dimensional-weak-strong-equivalence}，弱拓扑与强拓扑一致，故范数紧与弱紧等价。
\end{proof}

\subsection{存在性与内部概念的有限维简化}

\begin{theorem}[有限维塌缩原则]\label{thm:finite-dimensional-collapse-principle}
设 $X=\mathbb R^n$。则前七章中与存在性、对偶性和 KKT 相关的许多技术门槛都可以统一简化为有限维表述：
\begin{enumerate}
  \item 弱收敛、强收敛与坐标收敛一致；
  \item 闭有界可行集自动紧，因此第 \ref{sec:5-1} 节直接法中的弱紧性条件可压缩成闭有界性；
  \item 对非空凸集，拓扑内部、相对内部与第 \ref{sec:4-2} 节中的几何余量更容易协调；
  \item 第 \ref{sec:7-4} 节定义~\ref{def:kkt-cone-system} 的 KKT 系统在矩阵坐标下变成有限条代数方程与不等式。
\end{enumerate}
\end{theorem}

\begin{proof}
第一条由命题~\ref{prop:finite-dimensional-weak-strong-equivalence} 直接给出。第二条由推论~\ref{cor:finite-dimensional-compactness} 得到，而这正是第 \ref{sec:3-2} 节推论~\ref{cor:hilbert-weak-compactness} 在有限维里的强化版本。第三条的原因在于，有限维凸集的仿射包本身还是有限维欧氏空间，因此相对内部可直接用 Euclidean 邻域刻画，并与常用图形直觉保持一致。第四条则来自于：有限维对偶空间仍是有限维向量空间，连续线性算子都可由矩阵表示，故抽象伴随、乘子与互补松弛都能写成坐标化代数系统。
\end{proof}

\begin{remark}
定理~\ref{thm:finite-dimensional-collapse-principle} 的真正价值，是解释 Boyd 书中那些看起来“不言自明”的步骤到底省略了什么。譬如“取极小化列并抽收敛子列”“默认 KKT 乘子是向量”“相对内部和内部差别不大”等叙述，并不是普遍真理，而是有限维塌缩后的副产品。
\end{remark}

\subsection{典型例子}

\begin{example}[弱紧性与强紧性在有限维中一致]
在 $X=\mathbb R^n$ 中，任意闭球
\[
\{x\in \mathbb R^n:\|x\|_2\le r\}
\]
既是范数紧集，也是弱紧集。把这个例子放在这里，是为了把第 \ref{sec:3-2} 节中弱拓扑和强拓扑的技术分工一次性压缩回有限维直觉。
\end{example}

\begin{example}[相对内部与图形直觉]
对多面体
\[
C=\{x\in \mathbb R^n:Ax\le b\},
\]
相对内部可以直接理解为“在其仿射包里留有 Euclidean 余量的点”。把这个例子放在这里，是为了说明 Boyd 中大量关于相对内部和严格可行性的图形叙述，在有限维里是可靠的，但这种可靠性本身来自有限维塌缩。
\end{example}

\begin{example}[参数向量优化]
在机器学习和信号处理的标准模型里，参数往往直接属于 $\mathbb R^n$。把这个例子放在这里，是为了提醒读者：很多工程问题之所以能直接调用 Boyd 的有限维工具，不是因为它们先天地排斥无穷维理论，而是因为它们恰好处在定理~\ref{thm:finite-dimensional-collapse-principle} 覆盖的环境中。
\end{example}

\subsection*{有限维对照}

Boyd 的写法之所以可以把注意力几乎全部放在图像、矩阵和拉格朗日函数上，是因为弱拓扑、弱紧性和对偶空间的技术难点已经被有限维结构吸收了。本节的作用，就是把这种吸收过程显式讲清楚：抽象工具没有消失，它们只是退化成了 Euclidean 背景的默认设置。

\subsection*{习题（待补）}

\begin{exercise}[有限维塌缩占位]
补一题：用命题~\ref{prop:finite-dimensional-weak-strong-equivalence} 和推论~\ref{cor:finite-dimensional-compactness} 说明为什么有限维闭有界凸优化问题的极小值存在证明可以不用显式引入弱拓扑。
\end{exercise}

\begin{exercise}[相对内部桥接题占位]
补一题：对一个具体多面体写出其内部和相对内部，并指出在其仿射包降维后两者如何对应。
\end{exercise}


\section{线性规划与二次规划的复盘}

\label{sec:8-2}

第 \ref{sec:7-2} 节已经给出了抽象锥规划标准型，第 \ref{sec:7-4} 节则给出了无穷维 KKT 系统。本节要做的，不是重新讲一遍教科书里的 LP/QP 基础，而是把这些抽象对象逐项坍缩回矩阵语言。也正因此，线性规划和二次规划在这里的角色更像“坐标化译本”：同一个原--对偶结构，在有限维里被写成矩阵、向量和互补松弛式。

\subsection{线性规划与其对偶}

\begin{definition}[有限维线性规划]\label{def:linear-program-finite}
设 $c\in \mathbb R^n$，$A\in \mathbb R^{p\times n}$，$G\in \mathbb R^{m\times n}$，$b\in \mathbb R^p$，$h\in \mathbb R^m$。称问题
\[
\min_{x\in \mathbb R^n}\ c^\top x
\qquad \text{subject to}\qquad
Ax=b,\quad Gx\le h
\]
为一个\emph{线性规划}。
\end{definition}

\begin{proposition}[LP 的矩阵对偶形式]\label{prop:lp-dual-matrix-form}
定义~\ref{def:linear-program-finite} 的线性规划，其拉格朗日函数可写为
\[
L(x,\nu,\lambda)=c^\top x+\nu^\top(Ax-b)+\lambda^\top(Gx-h),
\qquad \lambda\ge 0.
\]
相应对偶问题为
\[
\max_{\nu\in \mathbb R^p,\ \lambda\in \mathbb R_+^m}\ -b^\top \nu-h^\top \lambda
\]
subject to
\[
c+A^\top \nu+G^\top \lambda=0.
\]
\end{proposition}

\begin{proof}
这正是第 \ref{sec:7-2} 节定义~\ref{def:abstract-cone-program} 在 $K=\mathbb R_+^m$ 且等式约束单独列出的有限维形式。对偶函数定义为
\[
q(\nu,\lambda):=\inf_{x\in \mathbb R^n}L(x,\nu,\lambda).
\]
若
\[
c+A^\top \nu+G^\top \lambda\neq 0,
\]
则 $L(x,\nu,\lambda)$ 对 $x$ 是非常数仿射函数，其下确界为 $-\infty$。只有当
\[
c+A^\top \nu+G^\top \lambda=0
\]
时，对偶函数才有限，此时
\[
q(\nu,\lambda)=-b^\top \nu-h^\top \lambda.
\]
再结合 $\lambda\ge 0$，便得到所述矩阵对偶形式。
\end{proof}

\begin{theorem}[LP 的矩阵型 KKT 条件]\label{thm:lp-kkt-matrix}
设线性规划存在最优解，并满足有限维 Slater 型严格可行性。则 $x^\star$ 为原问题最优解、$(\nu^\star,\lambda^\star)$ 为对偶最优解，当且仅当
\[
Ax^\star=b,\qquad Gx^\star\le h,
\]
\[
\lambda^\star\ge 0,
\]
\[
c+A^\top \nu^\star+G^\top \lambda^\star=0,
\]
\[
\lambda_i^\star(h_i-g_i^\top x^\star)=0,\qquad i=1,\dots,m.
\]
\end{theorem}

\begin{proof}
这正是第 \ref{sec:7-4} 节定理~\ref{thm:infinite-dimensional-kkt} 在有限维 Euclidean 空间中的坍缩。原可行性与对偶可行性分别对应前两组和第三组条件；驻点条件由
\[
0\in \partial(c^\top x)+A^\top \nu^\star+G^\top \lambda^\star
\]
简化为线性方程
\[
c+A^\top \nu^\star+G^\top \lambda^\star=0;
\]
互补松弛则由
\[
\langle \lambda^\star,h-Gx^\star\rangle=0
\]
在逐项非负锥上的配对写成逐坐标乘积等于零。
\end{proof}

\subsection{二次规划}

\begin{definition}[有限维二次规划]\label{def:quadratic-program-finite}
设 $H\in \mathbb S^n_+$ 为对称半正定矩阵，$c\in \mathbb R^n$。称问题
\[
\min_{x\in \mathbb R^n}\ \frac12 x^\top Hx+c^\top x
\qquad \text{subject to}\qquad
Ax=b,\quad Gx\le h
\]
为一个\emph{凸二次规划}。
\end{definition}

\begin{proposition}[QP 的矩阵型 KKT 系统]\label{prop:qp-kkt-matrix}
在定义~\ref{def:quadratic-program-finite} 中，若满足有限维严格可行性，则 $x^\star$ 为最优解，当且仅当存在 $(\nu^\star,\lambda^\star)$ 使
\[
Ax^\star=b,\qquad Gx^\star\le h,\qquad \lambda^\star\ge 0,
\]
\[
Hx^\star+c+A^\top \nu^\star+G^\top \lambda^\star=0,
\]
\[
\lambda_i^\star(h_i-g_i^\top x^\star)=0,\qquad i=1,\dots,m.
\]
\end{proposition}

\begin{proof}
由 $H\succeq 0$，目标函数是凸且可微，其梯度为
\[
\nabla\!\left(\frac12 x^\top Hx+c^\top x\right)=Hx+c.
\]
因此第 \ref{sec:5-2} 节的一阶微分语言和第 \ref{sec:7-4} 节的 KKT 系统结合后，驻点条件正好化为
\[
Hx^\star+c+A^\top \nu^\star+G^\top \lambda^\star=0.
\]
其余条件与线性规划完全一致，故结论成立。
\end{proof}

\begin{remark}
命题~\ref{prop:qp-kkt-matrix} 说明，第 \ref{sec:7-4} 节中的抽象乘子系统在有限维里会同时吸收第 \ref{sec:5-2} 节的梯度与第 \ref{sec:5-3} 节的次微分语言。LP 对应的是线性目标的次梯度为常向量，QP 对应的是二次目标的梯度为仿射映射；而结构本身并没有变化。
\end{remark}

\subsection{典型例子}

\begin{example}[多面体上线性目标]
线性规划中的可行集是多面体，目标则是线性泛函。把这个例子放在这里，是为了把第 \ref{sec:4-4} 节极点与支撑超平面的语言回收到最经典的 LP 场景：抽象分离与极点理论，在有限维里正好对应多面体几何。
\end{example}

\begin{example}[标准二次规划]
当 $H$ 为对称半正定矩阵时，定义~\ref{def:quadratic-program-finite} 给出的 QP 同时体现了第 \ref{sec:5-2} 节的梯度、第 \ref{sec:7-2} 节的标准型和第 \ref{sec:7-4} 节的 KKT 条件。把这个例子放在这里，是为了让读者看到 Hessian、乘子和互补松弛在有限维中如何自然拼接起来。
\end{example}

\begin{example}[均值--方差投资组合]
均值--方差组合优化可写为
\[
\min_x \frac12 x^\top \Sigma x-\mu^\top x
\qquad \text{subject to}\qquad
\mathbf 1^\top x=1,\quad x\ge 0.
\]
把这个例子放在这里，是为了给第 \ref{chap:11} 章连续时间模型一个静态有限维基线，同时说明 QP 的矩阵型 KKT 公式具有直接的经济解释。
\end{example}

\subsection*{有限维对照}

Boyd 书中 LP 和 QP 的拉格朗日对偶与 KKT 条件通常直接写成矩阵形式，这会给人一种错觉：似乎这些结论只属于线性代数世界。本节的作用正相反。通过回链定义~\ref{def:abstract-cone-program}、定理~\ref{thm:infinite-dimensional-kkt} 与推论~\ref{cor:kkt-via-fenchel-subdifferential}，我们看到矩阵公式只是抽象 KKT 在 $\mathbb R^n$ 中的坐标表示。

\subsection*{习题（待补）}

\begin{exercise}[LP 对偶占位]
补一题：从定义~\ref{def:linear-program-finite} 出发，直接计算线性规划的对偶函数并推出命题~\ref{prop:lp-dual-matrix-form}。
\end{exercise}

\begin{exercise}[QP KKT 桥接题占位]
补一题：对一个带不等式约束的凸二次规划写出命题~\ref{prop:qp-kkt-matrix} 的全部条件，并说明哪些部分来自第 \ref{sec:7-4} 节的抽象 KKT 系统。
\end{exercise}


\section{SOCP 与 SDP}

\label{sec:8-3}

如果说上一节把一般锥规划退回到了 LP/QP，那么本节要展示的则是最具代表性的两个有限维锥：二阶锥和半正定锥。它们的重要性不只在于应用面广，更在于它们恰好把第 \ref{sec:7-1} 节的广义不等式语言和第 \ref{sec:3-4} 节的谱论语言接到了一起。SOCP 强调“锥”的几何自对偶性，SDP 则强调“正性”与“谱分解”的矩阵化。

\subsection{二阶锥与半正定锥}

\begin{definition}[二阶锥]\label{def:second-order-cone}
在 $\mathbb R\times \mathbb R^{m-1}$ 中，定义
\[
\mathcal Q_m:=\{(t,z):t\ge \|z\|_2\}.
\]
该集合称为 $m$ 维\emph{二阶锥}，也称 Lorentz 锥。
\end{definition}

\begin{definition}[半正定锥]\label{def:positive-semidefinite-cone}
记 $\mathbb S^n$ 为 $n\times n$ 实对称矩阵空间。定义
\[
\mathbb S_+^n:=\{X\in \mathbb S^n:x^\top Xx\ge 0,\ \forall x\in \mathbb R^n\}.
\]
称其为\emph{半正定锥}。
\end{definition}

\begin{proposition}[二阶锥的自对偶性]\label{prop:soc-self-dual}
在标准 Euclidean 内积下，二阶锥满足
\[
\mathcal Q_m^\ast=\mathcal Q_m.
\]
\end{proposition}

\begin{proof}
设 $(s,y)\in \mathcal Q_m$。对任意 $(t,z)\in \mathcal Q_m$，由 Cauchy--Schwarz 不等式，
\[
s t+y^\top z\ge st-\|y\|_2\|z\|_2\ge st-\|y\|_2 t=t(s-\|y\|_2)\ge 0,
\]
故 $\mathcal Q_m\subset \mathcal Q_m^\ast$。

反之，若 $(s,y)\in \mathcal Q_m^\ast$，则对所有 $(t,z)\in \mathcal Q_m$ 有
\[
st+y^\top z\ge 0.
\]
特别取 $t=\|z\|_2$ 且令 $z=-\|z\|_2\,y/\|y\|_2$（当 $y\neq 0$ 时），便得
\[
\|z\|_2(s-\|y\|_2)\ge 0.
\]
由 $\|z\|_2>0$ 任意，可知 $s\ge \|y\|_2$，故 $(s,y)\in \mathcal Q_m$。于是 $\mathcal Q_m^\ast=\mathcal Q_m$。
\end{proof}

\begin{proposition}[迹内积与半正定锥的对偶配对]\label{prop:trace-inner-product-duality}
在 $\mathbb S^n$ 上定义迹内积
\[
\langle X,Y\rangle:=\operatorname{tr}(XY).
\]
则对任意 $X,Y\in \mathbb S_+^n$，
\[
\langle X,Y\rangle\ge 0,
\]
并且
\[
(\mathbb S_+^n)^\ast=\mathbb S_+^n
\]
其中对偶是相对于迹内积而言。
\end{proposition}

\begin{proof}
若 $X,Y\succeq 0$，则存在对称平方根 $Y^{1/2}$，从而
\[
\operatorname{tr}(XY)=\operatorname{tr}(Y^{1/2}XY^{1/2})\ge 0
\]
因为矩阵 $Y^{1/2}XY^{1/2}$ 仍半正定，其特征值非负，迹自然非负。

反之，若某个对称矩阵 $Z$ 满足
\[
\operatorname{tr}(ZY)\ge 0,\qquad \forall Y\succeq 0,
\]
则特别对秩一矩阵 $Y=uu^\top\succeq 0$ 有
\[
u^\top Zu=\operatorname{tr}(Zuu^\top)\ge 0,\qquad \forall u\in \mathbb R^n.
\]
故 $Z\succeq 0$。这就证明了半正定锥对迹内积是自对偶的，而迹内积正是第 \ref{sec:7-1} 节定义~\ref{def:dual-cone} 在矩阵空间中的具体写法。
\end{proof}

\subsection{谱语言在矩阵上的坍缩}

\begin{theorem}[对称矩阵的谱分解坍缩]\label{thm:spectral-collapse-symmetric-matrix}
设 $X\in \mathbb S^n$。则存在正交矩阵 $Q$ 与实对角矩阵
\[
\Lambda=\operatorname{diag}(\lambda_1,\dots,\lambda_n)
\]
使
\[
X=Q\Lambda Q^\top.
\]
特别地，
\[
X\succeq 0
\quad\Longleftrightarrow\quad
\lambda_i\ge 0,\ i=1,\dots,n.
\]
\end{theorem}

\begin{proof}
这正是第 \ref{sec:3-4} 节定理~\ref{thm:compact-selfadjoint-spectral} 在有限维 Hilbert 空间 $\mathbb R^n$ 上的坍缩。有限维下所有线性算子都是紧算子，因此任意自伴算子都可正交对角化。对 $X\in \mathbb S^n$ 应用该谱定理，即得
\[
X=Q\Lambda Q^\top.
\]
再注意到对任意 $u\in \mathbb R^n$，令 $v=Q^\top u$，则
\[
u^\top Xu=v^\top \Lambda v=\sum_{i=1}^n \lambda_i v_i^2.
\]
因此 $u^\top Xu\ge 0$ 对所有 $u$ 成立，当且仅当每个 $\lambda_i\ge 0$。
\end{proof}

\begin{remark}
定理~\ref{thm:spectral-collapse-symmetric-matrix} 解释了 SDP 为什么同时属于锥优化和谱论两条主线：矩阵正性既可以理解为“落在半正定锥中”，也可以理解为“所有特征值非负”。这两种说法在有限维中完全等价，而在无穷维里则必须借助第 \ref{sec:3-4} 节的紧自伴谱定理才能恢复。
\end{remark}

\subsection{典型例子}

\begin{example}[二阶锥几何]
二阶锥约束
\[
\|Fx+g\|_2\le a^\top x+b
\]
可直接写成
\[
(a^\top x+b,\,Fx+g)\in \mathcal Q_m.
\]
把这个例子放在这里，是为了说明 SOCP 并不是“带范数的不等式技巧”，而是第 \ref{sec:7-1} 节广义不等式语言在一个自对偶锥上的直接落地。
\end{example}

\begin{example}[对称矩阵特征值分解]
对任意 $X\in \mathbb S^n$，定理~\ref{thm:spectral-collapse-symmetric-matrix} 给出的
\[
X=Q\Lambda Q^\top
\]
就是有限维最熟悉的特征值分解。把这个例子放在这里，是为了把 Boyd 关于 SDP 的几何图像、矩阵迹内积和半正定约束统一到一个谱分解视角中。
\end{example}

\begin{example}[LMI 与协方差估计]
鲁棒控制和协方差估计中常见约束是
\[
A_0+\sum_{i=1}^n x_i A_i\succeq 0,
\]
这是一类线性矩阵不等式。把这个例子放在这里，是为了给工程读者一个直接触点：LMI 只是半正定锥约束的 affine 切片，因此它们天然属于第 \ref{sec:7-1} 节的锥规划框架。
\end{example}

\subsection*{有限维对照}

Boyd 书中 SOCP 和 SDP 往往被当成“高级有限维模型”单独介绍，但从本书视角看，它们不过是两件事情的交叉：第 \ref{sec:7-1} 节的锥语言，以及第 \ref{sec:3-4} 节的谱论语言。有限维之所以让这些模型显得格外整洁，是因为对偶锥、自伴正性和特征值分解都能同时写成矩阵公式。

\subsection*{习题（待补）}

\begin{exercise}[SOCP 自对偶占位]
补一题：直接验证命题~\ref{prop:soc-self-dual}，并把二阶锥约束改写为一个广义不等式。
\end{exercise}

\begin{exercise}[SDP 谱分解桥接题占位]
补一题：利用定理~\ref{thm:spectral-collapse-symmetric-matrix} 证明矩阵半正定等价于所有特征值非负，并说明这与命题~\ref{prop:trace-inner-product-duality} 的关系。
\end{exercise}


\section{几何规划的对偶透视}

\label{sec:8-4}

几何规划表面上似乎离前几章很远，因为它天然带有乘法结构、幂函数和 posynomial 目标，而不是线性算子与加法对偶。但这正是它适合放在本章末尾的原因：GP 是一个最典型的“特例坍缩”。一旦做对数变换，乘法会变成加法，monomial 会变成仿射函数，posynomial 会变成 log-sum-exp，从而整个问题重新回到凸分析和 Fenchel--Rockafellar 对偶的轨道上。

\subsection{几何规划的基本形式}

\begin{definition}[几何规划]\label{def:geometric-program}
一个标准\emph{几何规划}通常写成
\[
\min_{x\in \mathbb R_{++}^n}\ f_0(x)
\]
subject to
\[
f_i(x)\le 1,\qquad i=1,\dots,m,
\]
\[
g_j(x)=1,\qquad j=1,\dots,p,
\]
其中每个 $f_i$ 是 posynomial，每个 $g_j$ 是 monomial。这里
\[
g(x)=c x_1^{a_1}\cdots x_n^{a_n},\qquad c>0,
\]
称为 monomial；若
\[
f(x)=\sum_{k=1}^N c_k x_1^{a_{1k}}\cdots x_n^{a_{nk}},\qquad c_k>0,
\]
则称为 posynomial。
\end{definition}

\begin{proposition}[对数变换后的凸性]\label{prop:gp-log-transform-convexity}
令
\[
y_i=\log x_i,\qquad x_i=e^{y_i}.
\]
则 monomial 的对数变成仿射函数，posynomial 的对数变成 convex 的 log-sum-exp 形式。特别地，几何规划在对数坐标下可转化为凸优化问题。
\end{proposition}

\begin{proof}
对 monomial，
\[
\log g(e^y)=\log c+\sum_{i=1}^n a_i y_i,
\]
这是仿射函数。

对 posynomial，
\[
f(e^y)=\sum_{k=1}^N c_k \exp(a_k^\top y),
\]
故
\[
\log f(e^y)=\log \left(\sum_{k=1}^N \exp(a_k^\top y+\log c_k)\right).
\]
右端正是经典的 log-sum-exp 形式。由于指数函数凸且求和保持凸性，而对数作用在正和式上得到的 log-sum-exp 仍是凸函数，因此在变量 $y$ 上得到的是凸约束与凸目标。
\end{proof}

\begin{definition}[GP 中的 log-sum-exp 目标]\label{def:log-sum-exp-gp}
若 $f$ 是 posynomial，则在对数坐标 $y=\log x$ 下定义
\[
F(y):=\log f(e^y)
=
\log \left(\sum_{k=1}^N \exp(a_k^\top y+\beta_k)\right),
\qquad \beta_k:=\log c_k.
\]
函数 $F$ 称为 GP 变换后的\emph{log-sum-exp} 目标或约束函数。
\end{definition}

\subsection{GP 的对偶透视}

\begin{theorem}[GP 的 Fenchel 对偶表示]\label{thm:gp-duality-fenchel}
在对数变量 $y$ 下，几何规划可写成一个由仿射项与 log-sum-exp 项组成的有限维凸优化问题。因而在满足有限维资格条件时，其对偶问题可由第 \ref{sec:6-4} 节定理~\ref{thm:fenchel-rockafellar-duality} 给出；换言之，GP 的经典对偶本质上是 log-sum-exp 共轭与仿射约束共同作用的 Fenchel 对偶。
\end{theorem}

\begin{proof}
由命题~\ref{prop:gp-log-transform-convexity}，原 GP 在对数变量下化为
\[
\min_y\ F_0(y)
\]
subject to
\[
F_i(y)\le 0,\qquad i=1,\dots,m,
\]
\[
\alpha_j^\top y+\gamma_j=0,\qquad j=1,\dots,p,
\]
其中每个 $F_i$ 都是定义~\ref{def:log-sum-exp-gp} 给出的 log-sum-exp 函数。于是这已经是一个标准有限维凸优化问题。

再把不等式约束写成指示函数，把等式约束写成仿射子空间的指示函数，便可将整个问题整理为
\[
\inf_y \{\Phi(y)+\Psi(By)\},
\]
其中 $\Phi$ 由 log-sum-exp 项组成，$\Psi$ 由约束指示函数组成。此时第 \ref{sec:6-1} 节定义~\ref{def:fenchel-conjugate} 的 Fenchel 共轭和第 \ref{sec:6-4} 节定理~\ref{thm:fenchel-rockafellar-duality} 直接适用，于是得到 GP 的对偶表示。也就是说，GP 并不是对偶理论之外的特例，而是对数坐标下的有限维对偶坍缩。
\end{proof}

\begin{remark}
从本书视角看，GP 最值得强调的不是“它可以变凸”这一口号，而是“它变凸之后，所依赖的仍然是同一套共轭与对偶框架”。这就是为什么本节要显式回链第 \ref{sec:6-4} 节，而不是把 GP 当作一组孤立技巧。
\end{remark}

\subsection{典型例子}

\begin{example}[功率控制的几何规划]\label{ex:power-control-gp}
无线功率控制中的信噪比约束常可整理为
\[
\frac{\sigma_i+\sum_{j\neq i} a_{ij}p_j}{p_i}\le \frac{1}{\gamma_i},
\qquad p_i>0.
\]
在正变量 $p_i$ 上这看起来是分式与乘法混合结构；经对数变换后，它会变成 log-sum-exp 约束。把这个例子放在这里，是为了回收 Boyd 中最经典的 GP 工程模型，并说明它本质上属于有限维对偶坍缩。
\end{example}

\begin{example}[电路与面积功率折中]
集成电路设计里经常出现 monomial 与 posynomial 形式的面积、延迟和功耗模型。把这个例子放在这里，是为了说明 GP 不是数学上的边角料，而是对数凸建模在工程中的代表性落点。
\end{example}

\begin{example}[乘法效用资源分配]
若效用函数带有乘法结构，如
\[
\prod_{i=1}^n x_i^{\alpha_i},
\]
则取对数后会转化为加法目标
\[
\sum_{i=1}^n \alpha_i \log x_i.
\]
把这个例子放在这里，是为了给管理科学读者一个直观触点：乘法效用模型之所以适合凸优化，不是因为乘法本身“神奇”，而是因为对数坐标把它送回了加法与对偶的世界。
\end{example}

\subsection*{有限维对照}

Boyd 书中几何规划往往以对数变换和工程例子同时出现，因此读者容易把它记成一套技巧性很强的特例。本节强调的恰恰是相反的事实：GP 在有限维里之所以漂亮，是因为 log-sum-exp、Fenchel 共轭和 Fenchel--Rockafellar 对偶都能被完整显式写出。也就是说，GP 不是对凸对偶的例外，而是它最精致的坍缩之一。

\subsection*{习题（待补）}

\begin{exercise}[GP 对数变换占位]
补一题：把一个具体的 posynomial 目标和约束问题改写到对数变量下，并验证 log-sum-exp 函数的出现。
\end{exercise}

\begin{exercise}[GP 对偶桥接题占位]
补一题：对一个简单的几何规划写出其对数凸化问题，并说明第 \ref{sec:6-4} 节 Fenchel--Rockafellar 对偶如何给出对应对偶表达。
\end{exercise}

